\chapter{From convincing examples to proof}
\label{learn:proof}

\begin{learningcard}{Your starting equipment}
Bring integer arithmetic, algebraic identities, inequalities, functions, and
finite sums from the precalculus book. The new skill is to specify exactly
which objects a statement covers and why the reasoning covers every one of
them. The question guiding the chapter is: when does a pattern become a theorem?
\end{learningcard}

\section{A statement needs a universe}
A proposition is a statement with a truth value once its terms and domain are
fixed. ``Every number has a square root'' changes meaning when ``number''
means a nonnegative real, an arbitrary real, or a complex number. We write
$\RR$ for the real numbers and $\ZZ$ for the integers. The symbols
$\forall$ and $\exists$ mean ``for every'' and ``there exists.'' Thus
\[
 \forall x\in\RR\;\exists y\in\RR:\quad y>x
\]
says that each real number has a larger real number. The choice $y=x+1$
proves the assertion. Reversing the quantifiers asks for one real number
larger than every real number. Any proposed $y$ fails at $x=y+1$.
The order records which choices may depend on earlier choices.

An implication $P\Rightarrow Q$ says that whenever the hypothesis $P$ holds,
the conclusion $Q$ holds. Its converse $Q\Rightarrow P$ is a separate claim.
For integers, divisibility by four implies evenness; the even integer six
refutes the converse. An equivalence $P\Leftrightarrow Q$ requires both
implications. A theorem is a proved proposition within stated assumptions;
a conjecture is a proposition offered for investigation.

\begin{learningcard}{Historical situation: a diagram and its obligations}
Euclid's \emph{Elements}, Book I, Proposition 47 proves the square relation
for a right triangle through constructions and earlier propositions. The
surviving text gives a structured chain of dependencies rather than a table
of measured triangles. Its transmission, translation, and editorial history
are separate historical questions. We use the proposition as evidence of
an organized deductive presentation, rather than as evidence that one person
invented every result in the collection. David Joyce's mathematical edition
provides the proposition, diagram, and commentary.\footnote{Euclid,
\emph{Elements}, I.47, edition and commentary by David E. Joyce, Clark
University: \url{https://mathcs.clarku.edu/~djoyce/elements/bookI/propI47.html}.}
The modern logical notation in this chapter is our teaching apparatus.
\end{learningcard}

\section{Worked example: move from a definition to a conclusion}
An integer $m$ is even when $m=2a$ for some $a\in\ZZ$, and odd when
$m=2a+1$ for some $a\in\ZZ$. Prove that the product of two odd integers
is odd. Start with arbitrary odd integers, writing $m=2a+1$, $n=2b+1$.
Then
\[
 mn=(2a+1)(2b+1)=4ab+2a+2b+1=2(2ab+a+b)+1.
\]
The expression $2ab+a+b$ is an integer because integers are closed under
addition and multiplication. The final expression therefore satisfies the
definition of oddness. The proof's decisive move is the extraction of an
integer witness, rather than the expansion alone.

A proof by contrapositive establishes $P\Rightarrow Q$ by proving
$\neg Q\Rightarrow\neg P$. For example, if $m^2$ is even, then $m$ is even:
if $m$ were odd, the product result would make $m^2$ odd. Every integer
is either even or odd, which completes the implication. A proof by
contradiction instead assumes the desired conclusion false together with
the hypotheses and derives an impossibility. These are related logical
strategies, but the argument should state which assumption it introduces.

\section{Worked example: induction is a dependency chain}
To prove a statement $P(n)$ for every integer $n\geq1$, mathematical
induction requires a base case $P(1)$ and an implication
$P(k)\Rightarrow P(k+1)$ for every integer $k\geq1$. The implication
propagates truth from the base through the positive integers.
Consider
\[
 1+2+\cdots+n=\frac{n(n+1)}2.
\]
At $n=1$, both sides equal one. Assume the formula for an arbitrary
$k\geq1$. Add the next term to obtain
\[
 \sum_{j=1}^{k+1}j=\frac{k(k+1)}2+(k+1)
 =\frac{(k+1)(k+2)}2.
\]
The right side is precisely the formula at $k+1$. The assumption covers
one unspecified preceding case; induction then supplies all cases.
Checking the first thousand cases alone would leave every later integer
outside the checked set.

The base case carries real information. The implication
``if $2^k$ is odd, then $2^{k+1}$ is odd'' holds for $k\geq1$ because
its hypothesis always fails, yet $2^n$ is even for every positive $n$.
A chain of implications needs a true starting point.

\section{A counterexample has a small job}
A universal claim fails when one admissible object violates its conclusion.
For ``every prime is odd,'' two is sufficient: it is prime and even.
A counterexample must satisfy the hypotheses. The number one fails as a
counterexample because a prime is an integer greater than one whose only
positive divisors are one and itself.

An existential statement needs one witness. A universal statement over a
finite, completely enumerated domain can be proved by exact exhaustive
calculation. The same calculation over a sample from an infinite domain
establishes only the sampled cases. Floating-point computations introduce
a further boundary: a small residual measures agreement to a tolerance,
whereas exact equality requires an exact argument or exact arithmetic.

\begin{learningcard}{Choose a useful next question}
Before calculating, write the domain, quantifiers, and assumptions. After
calculating, identify the witness or implication that makes the result
apply to its claimed domain. Continue to Chapter~\ref{learn:linear_algebra}
for structural proofs or Chapter~\ref{learn:calculus} for quantified limits.
\end{learningcard}

\section{Exercises with full solutions}
\paragraph{1. Negate a quantified statement.}
Negate: ``For every real $x$ there exists a real $y$ such that $y^2=x$.''
Determine the truth of the original and its negation.

\paragraph{Solution.}
The negation is ``There exists a real $x$ such that for every real $y$,
$y^2\ne x$.'' Choose $x=-1$. Every real square is nonnegative, so every
real $y$ satisfies $y^2\ne-1$. The negation is true and the original is
false. The witness refutes a universal statement; checking positive
values of $x$ would leave the negative part of the domain unresolved.

\paragraph{2. Prove an inequality by induction.}
Prove $2^n\geq n+1$ for every integer $n\geq0$.

\paragraph{Solution.}
At zero, $2^0=1=0+1$. Assuming $2^k\geq k+1$ for an arbitrary
$k\geq0$, multiplication by the positive number two gives
$2^{k+1}\geq2k+2$. Since $2k+2-(k+2)=k\geq0$, we obtain
$2^{k+1}\geq k+2$. Induction completes the proof, including the stated
zero endpoint.

\paragraph{3. Separate an implication and its converse.}
For a real number $x$, analyze ``$x>2$ implies $x^2>4$'' and its converse.

\paragraph{Solution.}
If $x>2$, then both $x-2$ and $x+2$ are positive, so
$x^2-4=(x-2)(x+2)>0$. The converse fails at $x=-3$, whose square is
nine. A corrected equivalence is $x^2>4$ if and only if $|x|>2$:
the factors show the quadratic is positive on $(-\infty,-2)$ and
$(2,\infty)$, exactly the domain described by the absolute-value inequality.
